\section{Moltiplicazione di matrici - Strassen}
\subsubsection*{Note introduttive sulle matrici}
Per la seguente sezione relativa al prodotto di matrici, si considerano solo matrici quadrate \(n \times n\)
(\(\in \mathbb{R}^{n \times n}\)) con \(n = 2^k\) potenza di 2, per un certo \(k \in \mathbb{N}\). In questo modo
le matrici sono perfettamente scomponibili ricorsivamente in quattro sottomatrici di dimensione \(n/2 \times n/2\).

Per convenzione si assume che la taglia di una matrice è pari al numero di righe (o di colonne) di cui è composta,
ovvero \(A \in \mathbb{R}^{n \times n} \;\Rightarrow\; |A| = n\).

\subsection{Prodotto di matrici usando la definizione}
\subsubsection*{Tecnica risolutiva}
Date due matrici \(A, B \in \mathbb{R}^{n \times n}\), con \(n = 2^k\), il loro prodotto \(C = A \cdot B\) è sempre una
matrice \(C \in \mathbb{R}^{n \times n}\) i cui elementi \(c_{ij}\) sono calcolati come segue:
\[c_{ij} = \sum_{k=1}^{n} a_{ik} \cdot b_{kj} \qquad\qquad \text{per } i, j = 1, 2, \dotsc, n\]
In altre parole il prodotto di matrici è una matrice che ha come elementi il prodotto interno dei vettori riga e colonna
rispettivamente della prima matrice \(A\) e della seconda matrice \(B\).

\subsubsection*{Analisi della complessità}
Il prodotto interno tra un vettore riga e un vettore colonna richiede \(n\) moltiplicazioni e \(n-1\) addizioni, per un costo
totale di \(T(n) = 2n - 1\). Siccome vanno fatti \(n^2\) prodotti interni, il costo totale della moltiplicazione di matrici
risulta \(\Theta(n^3)\): \[T(n) = n^2 (2n - 1) = 2n^3 - n^2 = \Theta(n^3)\]

\subsection{Prodotto di matrici con approccio divide and conquer}
\subsubsection*{Tecnica risolutiva}
Le matrici \(A, B \in \mathbb{R}^{n \times n}\) possono essere scomposte ciascuna in quattro sottomatrici di dimensioni
\(n/2 \times n/2\), identificate secondo il seguente sistema di indici \(X_{11}, X_{12}, X_{21}, X_{22} \in \mathbb{R}^{n/2 \times n/2}\).
Il prodotto \(C = A \cdot B\) risulta essere in funzione dei prodotti tra queste sottomatrici.
\[A = \begin{bmatrix}
	A_{11} & A_{12} \\ A_{21} & A_{22}
\end{bmatrix} \qquad\qquad\quad B = \begin{bmatrix}
	B_{11} & B_{12} \\ B_{21} & B_{22}
\end{bmatrix} \qquad\qquad\quad C = \begin{bmatrix}
	C_{11} & C_{12} \\ C_{21} & C_{22}
\end{bmatrix}\]
\[\begin{array}{c c}
	C_{11} = A_{11} \cdot B_{11} + A_{12} \cdot B_{21} \qquad\quad & \qquad\quad C_{21} = A_{21} \cdot B_{11} + A_{22} \cdot B_{21} \\[5pt]
	C_{12} = A_{11} \cdot B_{12} + A_{12} \cdot B_{22} \qquad\quad & \qquad\quad C_{22} = A_{21} \cdot B_{12} + A_{22} \cdot B_{22}
\end{array}\]

\subsubsection*{Equazione alle ricorrenze}
L'equazione alle ricorrenze per questo approccio è la seguente:
\[T(n) \; = \; \begin{cases}
	\; 1 & \text{se } n = 1 \\[3pt]
	\; 8 \cdot T(n/2) + 4 n^2 & \text{se } n > 1
\end{cases}\]
\begin{itemize}
	\item il caso base è rappresentato dal prodotto di due matrici \(1 \times 1\) che richiede un solo prodotto
	\item il passo ricorsivo prevede 8 chiamate ricorsive per calcolare gli otto prodotti tra le sottomatrici di dimensione \(n/2 \times n/2\)
	\item il costo per ogni nodo interno è dato dalle 4 somme dei prodotti tra sottomatrici \(n/2 \times n/2\)
\end{itemize}

\subsubsection*{Analisi della complessità}
Dal formulone si ottiene che la complessità rimane sempre \(\Theta(n^3)\), come evidenziato di seguito:
\begin{align*}
	T(n) \;\; &= \sum_{l=0}^{(\log_2 n) - 1} \!\! \left( 8^l \cdot \left( \frac{n}{2^l} \right)^2 \right) + 8^{\log_2 n} \cdot T_0(n) \; = \; n^2 \cdot \! \sum_{l=0}^{(\log_2 n) - 1} \!\! 2^l + n^{\log_2 8} = n^2 \cdot \frac{2^{\log_2 n} - 1}{2 - 1} + n^3 \; = \\[3pt]
	&= \underbrace{\;n^2 \cdot (n - 1)\;}_{\begin{array}{c} \text{\scriptsize costo delle somme} \\[-3pt] \text{\scriptsize dei nodi interni} \end{array}}
	+ \!\!\!\!\!\! \underbrace{\;\;\; n^3 \;\;}_{\begin{array}{c} \text{\scriptsize costo dei prodotti} \\[-3pt] \text{\scriptsize alle foglie} \end{array}}
	\!\!\!\!\!\! = \;\; 2n^3 - n^2 \; = \; \Theta(n^3)
\end{align*}

\subsubsection*{Vantaggi dell'algoritmo D\&C ricorsivo rispetto a quello iterativo basato sulla definizione}
Ad oggi le macchine moderne performano molto meglio con l'algoritmo D\&C ricorsivo rispetto a quello iterativo basato sulla
definizione in quanto si lavora con sottomatrici sempre più piccole fino a quando riescono ad essere totalmente memorizzate
all'interno della cache. In questo modo si riducono drasticamente i tempi di accesso alla memoria principale dovuti ai
continui cache miss quando si scorrono per intero righe e colonne per calcolare i prodotti interni.

\subsection{Prodotto di matrici con l'algoritmo di Strassen (1969)}
\subsubsection*{Tecnica risolutiva}
Ispirandosi a Karatsuba, Strassen osserva che è possibile ricavare gli 8 prodotti tra le sottomatrici \(A_{ij}\) e \(B_{ij}\)
riutilizzando solo 7 prodotti, calcolati su opportune somme e differenze di sottomatrici.
\vspace{-12pt}
\begin{center}
	\begin{minipage}{0.25\textwidth}
		\begin{align*}
			A_1 &= (A_{12} - A_{22}) \\
			A_2 &= (A_{11} + A_{22}) \\
			A_3 &= (A_{11} - A_{21}) \\
			A_4 &= (A_{11} + A_{12}) \\
			A_5 &= A_{11} \\
			A_6 &= A_{22} \\
			A_7 &= (A_{21} + A_{22})
		\end{align*}
	\end{minipage}
	\begin{minipage}{0.25\textwidth}
		\begin{align*}
			B_1 &= (B_{21} + B_{22}) \\
			B_2 &= (B_{11} + B_{22}) \\
			B_3 &= (B_{11} + B_{12}) \\
			B_4 &= B_{22} \\
			B_5 &= (B_{12} - B_{22}) \\
			B_6 &= (B_{21} - B_{11}) \\
			B_7 &= B_{11}
		\end{align*}
	\end{minipage}
	\begin{minipage}{0.4\textwidth}
		\begin{align*}
			P_i = A_i \cdot B_i \qquad \text{per } i = 1, 2, \dotsc, 7
		\end{align*}
		\vspace{-20pt}
		\begin{align*}
			C_{11} &= P_1 + P_2 - P_4 + P_6 \\
			C_{12} &= P_4 + P_5 \\
			C_{21} &= P_6 + P_7 \\
			C_{22} &= P_2 - P_3 + P_5 - P_7
		\end{align*}
	\end{minipage}
\end{center}

\subsubsection*{Osservazioni sui prodotti di Strassen}
\begin{itemize}
	\item se si sviluppano i conti per le quattro espressioni di \(C_{ij}\) si può notare che si ottengono esattamente
	le stesse espressioni dell'approccio D\&C standard
	\item esiste una simmetria tra le espressioni dei blocchi \(A_i\) e \(B_i\): tra \(A_1\) e \(B_5\), tra \(A_2\)
	e \(B_2\), tra \(-A_3\) e \(B_6\), tra \(A_4\) e \(B_3\), tra \(A_5\) e \(B_7\), tra \(A_6\) e \(B_4\), tra \(A_7\) e \(B_1\)
\end{itemize}

\subsubsection*{Equazione alle ricorrenze}
L'equazione alle ricorrenze per l'algoritmo di Strassen è la seguente:
\[T(n) \; = \; \begin{cases}
	\; 1 & \text{se } n = 1 \\[3pt]
	\; 7 \cdot T(n/2) + 18 \; n^2/4 & \text{se } n > 1
\end{cases}\]

\begin{itemize}
	\item il caso base è rappresentato dal prodotto di due matrici \(1 \times 1\) che richiede un solo prodotto
	\item il passo ricorsivo prevede 7 chiamate ricorsive per calcolare i sette prodotti \(P_i\) tra i blocchi \(A_i\) e \(B_i\)
	di dimensione \(n/2 \times n/2\)
	\item il costo della fase di divide consiste nelle 10 somme e differenze tra sottomatrici \(A_{ij}, B_{ij} \in \mathbb{R}^{n/2 \times n/2}\)
	per calcolare i blocchi \(A_i\) e \(B_i\) (complessivamente \(10 \cdot n^2/4\))
	\item il costo della fase di conquer consiste nelle 8 somme e differenze tra i prodotti \(P_i \in \mathbb{R}^{n/2 \times n/2}\)
	per calcolare i blocchi \(C_{ij}\) (complessivamente \(8 \cdot n^2/4\))
\end{itemize}

\subsubsection*{Analisi della complessità}
Applicando il formulone si ottiene che la complessità è \(\Theta(n^{\log_2 7}) \approx \Theta(n^{2.81})\), come mostrato di seguito:
\begin{align*}
	T(n) \;\; &= \sum_{l=0}^{(\log_2 n) - 1} \!\! \left( 7^l \cdot \frac{18}{4} \cdot \left( \frac{n}{2^l} \right)^2 \right) + 7^{\log_2 n} \cdot T_0(n) \; = \; \frac{9}{2} \; n^2 \cdot \! \sum_{l=0}^{(\log_2 n) - 1} \!\! \left( (7/4)^l \right) + n^{\log_2 7} \; = \\[3pt]
	&= \;\; \frac{9}{2} \; n^2 \cdot \frac{(7/4)^{\log_2 n} - 1}{(7/4) - 1} + n^{\log_2 7} \; = \; \frac{9}{2} n^2 \cdot \frac{n^{\log_2 (7/4)} - 1}{3/4} + n^{\log_2 7} \; = \\[3pt]
	&= \;\; \frac{9}{2} \cdot \frac{4}{3} \; n^2 \cdot \left( n^{(\log_2 7) - 2} - 1 \right) + n^{\log_2 7} \; = \; 6 n^2 \cdot (n^{\log_2 7}/n^2 - 1) + n^{\log_2 7} \\[3pt]
	&= \; \underbrace{\; 6 n^{\log_2 7} - 6n^2 \;}_{\begin{array}{c} \text{\scriptsize costo delle somme} \\[-3pt] \text{\scriptsize dei nodi interni} \end{array}}
	\;\; + \!\!\!\! \underbrace{\;\; n^{\log_2 7} \;\;}_{\begin{array}{c} \text{\scriptsize costo dei prodotti} \\[-3pt] \text{\scriptsize alle foglie} \end{array}}
	\!\!\!\!\!\! = \;\; 7n^{\log_2 7} - 6n^2 \; = \; \Theta(n^{\log_2 7}) \; \approx \; \Theta(n^{2.81})
\end{align*}

\subsubsection*{Osservazioni generali sull'algoritmo di Strassen}
\begin{itemize}
	\item una volta le moltiplicazioni di numeri erano molto costose, ora i blocchi logici sono altamente ottimizzati; inoltre
	molti processori dispongono dell'istruzione MADD (multiply and add) che consente di accumulare il risultato di una
	moltiplicazione all'interno di un registro in un solo ciclo di clock
	\item se si volesse eseguire il quadrato di una matrice \(A\) con l'algoritmo di Strassen, si risparmierebbero 5 somme
	e sottrazioni siccome i blocchi \(A_i\) e \(B_j\) corrispondenti sarebbero uguali
	\item in generale non si può avere corrispondenza tra i blocchi \(A_i\) e \(B_i\), perché se così fosse, l'algoritmo
	diventerebbe commutativo portando sicuramente ad un risultato sbagliato siccome la moltiplicazione di matrici non è
	commutativa
	\item utilizzando la definizione, la complessità asintotica del prodotto di matrici era \(n^\Omega\) con \(\Omega = 3\),
	con Strassen si è riusciti a ridurre il valore di \(\Omega\) a \(\log_2 7 \approx 2.81\); ad oggi si conoscono algoritmi
	per il prodotto di matrici con \(\Omega \approx 2.37\), ma risultano poco vantaggiosi in quanto hanno costanti asintotiche
	molto elevate che li rendono svantaggiosi per valori di \(n\) piccoli
\end{itemize}

\subsection{Algoritmo ibrido per il prodotto di matrici}
\subsubsection*{Idea implementativa}
Per valori di \(n \leq n_0\) si osserva che l'algoritmo standard (D\&C) rimane ancora più efficiente di quello di Strassen.
Si vuole determinare tale valore \(n_0\) e sviluppare un algoritmo tale che per \(n > n_0\) utilizzi l'algoritmo di Strassen,
mentre per \(n < n_0\) utilizzi l'algoritmo standard (D\&C).

\subsubsection*{Calcolo del parametro \(n_0\)}
Per calcolare il parametro \(n_0\), si cerca il valore che minimizza la complessità dell'algoritmo di Strassen imponendo
che per \(n \leq n_0\) venga utilizzato l'algoritmo standard (D\&C), ovvero:
\begin{itemize}
	\item i livelli dell'albero della ricorsione arrivano a \((\log_2 (n/n_0)) - 1\)
	\item il costo delle foglie è pari alla complessità dell'algoritmo standard (D\&C), ovvero \(T_0(n) = 2n^3 - n^2\)
\end{itemize}
\begin{align*}
	T(n, n_0) \;\; &= \sum_{l=0}^{(\log_2 (n/n_0)) - 1} \!\! \left( 7^l \cdot \frac{18}{4} \cdot \left( \frac{n}{2^l} \right)^2 \right) \; + \; 7^{\log_2 (n/n_0)} \cdot \left( 2n_0^3 - n_0^2 \right) \; = \; \dots \; = \\[3pt]
	&= \;\; n^{\log_2 7} \cdot \underbrace{\frac{2 n_0 + 5}{{n_0}^{(\log_2 7) - 2}}}_{g(n_0)} - 6n^2
\end{align*}
Si procede a minimizzare la funzione \(T(n, n_0)\) per trovare il valore ottimale di \(n_0\):
\vspace{-5pt}
\[\min_{n_0} \, T(n, n_0) \;\; \longrightarrow \;\; \min_{n_0} \, g(n_0) \;\; \longrightarrow \;\; \begin{array}{l} n_{0, \min} \approx 10.477 \\[3pt] g(n_{0, \min}) = 3.895 \end{array}\]

\subsubsection*{Considerazioni sul parametro \(n_0\)}
Dal punto di vista teorico, si sceglie \(n_0 = 8\) in quanto è la potenza di 2 più vicina a \(n_{0, \min}\). Riscrivendo
l'equazione alle ricorrenze e la complessità asintotica dell'algoritmo ibrido, si ottiene:
\[T(n, n_{0, \min}) \; = \; \begin{cases}
	\; 2n^3 - n^2 & \text{se } n \leq 8 \\[3pt]
	\; 7 \cdot T(n/2) + 18 \; n^2/4 & \text{se } n > 8
\end{cases} \quad = \quad \begin{cases}
	\; 2n^3 - n^2 & \text{se } n \leq 8 \\[3pt]
	\; 3.918 \cdot n^{\log_2 7} - 6n^2 & \text{se } n > 8
\end{cases}\]

Dal punto di vista pratico, non è detto che \(n_0 = 8\) sia il valore ottimale, in quanto il modello teorico non tiene conto
di overhead dovuti al cambio di algoritmo, alla gestione della cache e altri fattori specifici della macchina su cui viene
eseguito l'algoritmo.

Ad oggi si sta sviluppando una branca di ricerca chiamata \textbf{adaptive software} che punta ad aggiustare dinamicamente i parametri
degli algoritmi di algebra lineare (come \(n_0\)) attraverso analisi e test pratici effettuati sulla macchina stessa, al fine
di ottenere le migliori prestazioni possibili per quella determinata macchina in uso.
